
\documentclass{article}
%%%%%%%%%%%%%%%%%%%%%%%%%%%%%%%%%%%%%%%%%%%%%%%%%%%%%%%%%%%%%%%%%%%%%%%%%%%%%%%%%%%%%%%%%%%%%%%%%%%%%%%%%%%%%%%%%%%%%%%%%%%%%%%%%%%%%%%%%%%%%%%%%%%%%%%%%%%%%%%%%%%%%%%%%%%%%%%%%%%%%%%%%%%%%%%%%%%%%%%%%%%%%%%%%%%%%%%%%%%%%%%%%%%%%%%%%%%%%%%%%%%%%%%%%%%%
\usepackage{amsmath, amssymb}
\usepackage{geometry}

\setcounter{MaxMatrixCols}{10}
%TCIDATA{OutputFilter=LATEX.DLL}
%TCIDATA{Version=5.50.0.2960}
%TCIDATA{<META NAME="SaveForMode" CONTENT="1">}
%TCIDATA{BibliographyScheme=Manual}
%TCIDATA{LastRevised=Monday, March 10, 2025 07:28:17}
%TCIDATA{<META NAME="GraphicsSave" CONTENT="32">}

\geometry{a4paper, margin=1in}

\input{tcilatex}

\begin{document}

\title{Relationship Between Dirac Matrices and Pauli Matrices}
\author{}
\date{}
\maketitle

\section*{Pauli Matrices}

The Pauli matrices $\sigma_1, \sigma_2, \sigma_3$ are three $2 \times 2$
complex matrices defined as: 
\begin{equation*}
\sigma_1 = 
\begin{pmatrix}
0 & 1 \\ 
1 & 0%
\end{pmatrix}%
, \quad \sigma_2 = 
\begin{pmatrix}
0 & -i \\ 
i & 0%
\end{pmatrix}%
, \quad \sigma_3 = 
\begin{pmatrix}
1 & 0 \\ 
0 & -1%
\end{pmatrix}%
. 
\end{equation*}
They satisfy the following properties:

\begin{itemize}
\item They are \textbf{Hermitian} ($\sigma_i^\dagger = \sigma_i$), 

\item They are \textbf{traceless} ($\text{Tr}(\sigma_i) = 0$), 

\item They satisfy the \textbf{Clifford algebra relations}:  
\begin{equation*}
\{ \sigma_i, \sigma_j \} = 2 \delta_{ij} I,  
\end{equation*}
where $\delta_{ij}$ is the Kronecker delta, and $I$ is the $2 \times 2$
identity matrix.
\end{itemize}

The Pauli matrices are used to describe \textbf{spin-1/2 particles} in
non-relativistic quantum mechanics.

\section*{Dirac Matrices}

The Dirac matrices $\gamma_\mu$ (where $\mu = 0, 1, 2, 3$) are four $4
\times 4$ complex matrices that satisfy the \textbf{Clifford algebra
relations} for Minkowski spacetime: 
\begin{equation*}
\{ \gamma_\mu, \gamma_\nu \} = 2 \eta_{\mu\nu} I, 
\end{equation*}
where $\eta_{\mu\nu}$ is the Minkowski metric with signature $(+, -, -, -)$,
and $I$ is the $4 \times 4$ identity matrix.

In the \textbf{Dirac representation}, the Dirac matrices are constructed
using the Pauli matrices as follows: 
\begin{equation*}
\gamma^0 = 
\begin{pmatrix}
I & 0 \\ 
0 & -I%
\end{pmatrix}%
, \quad \gamma^i = 
\begin{pmatrix}
0 & \sigma^i \\ 
-\sigma^i & 0%
\end{pmatrix}%
, 
\end{equation*}
where:

\begin{itemize}
\item $I$ is the $2 \times 2$ identity matrix, 

\item $\sigma^i$ ($i = 1, 2, 3$) are the Pauli matrices.
\end{itemize}

\section*{Relationship Between Dirac and Pauli Matrices}

The Dirac matrices are related to the Pauli matrices in the following ways:

\subsection*{Block Structure}

\begin{itemize}
\item The Dirac matrices are $4 \times 4$ matrices constructed from $2
\times 2$ blocks involving the Pauli matrices. 

\item The \textbf{spatial components} $\gamma^i$ ($i = 1, 2, 3$) directly
incorporate the Pauli matrices $\sigma^i$. 

\item The \textbf{timelike component} $\gamma^0$ is diagonal and does not
explicitly involve the Pauli matrices, but it still relates to the $2 \times
2$ structure of the Pauli matrices.
\end{itemize}

\subsection*{Spinor Representation}

\begin{itemize}
\item The Pauli matrices act on \textbf{2-component spinors} (Weyl spinors)
in 3-dimensional space. 

\item The Dirac matrices act on \textbf{4-component spinors} (Dirac spinors)
in 4-dimensional spacetime. These 4-component spinors can be decomposed into
two 2-component Weyl spinors:  
\begin{equation*}
\psi = 
\begin{pmatrix}
\psi_L \\ 
\psi_R%
\end{pmatrix}%
,  
\end{equation*}
where $\psi_L$ and $\psi_R$ are left-handed and right-handed Weyl spinors,
respectively.
\end{itemize}

\subsection*{Non-Relativistic Limit}

\begin{itemize}
\item In the \textbf{non-relativistic limit}, the Dirac equation reduces to
the \textbf{Pauli equation}, which describes spin-1/2 particles in
non-relativistic quantum mechanics. The Pauli matrices explicitly appear in
this limit.
\end{itemize}

\section*{Mathematical Connection}

The relationship between the Dirac and Pauli matrices can be understood
mathematically as follows:

\begin{itemize}
\item The Pauli matrices generate the \textbf{real Clifford algebra} $\text{%
Cl}(3, 0)$, which corresponds to 3-dimensional Euclidean space. 

\item The Dirac matrices generate the \textbf{real Clifford algebra} $\text{%
Cl}(1, 3)$, which corresponds to 4-dimensional Minkowski spacetime. 

\item The Dirac matrices extend the Pauli matrices to include the timelike
dimension, allowing for a description of relativistic spinors.
\end{itemize}

\section*{Summary}

\begin{itemize}
\item The \textbf{Dirac matrices} are constructed using the \textbf{Pauli
matrices} as building blocks. 

\item The Pauli matrices describe spinors in 3-dimensional space, while the
Dirac matrices describe spinors in 4-dimensional spacetime. 

\item The Dirac matrices extend the Pauli matrices to include the timelike
dimension, enabling a relativistic description of spin-1/2 particles.
\end{itemize}

\end{document}
